% !TEX program = xelatex
\documentclass[11pt,a4paper,fontset=none]{ctexart}
\usepackage[a4paper,margin=2.2cm]{geometry}
\usepackage{fancyhdr}
\setlength{\headheight}{14pt}
\pagestyle{fancy}
\fancyhf{}
\fancyhead[L]{ebbackup 答辩逐字稿}
\fancyhead[R]{\vtag{0.9.5}}
\fancyfoot[C]{\thepage}

% 答辩演讲稿 — 共享导言（颜色、盒子、tikz、listings）
\usepackage{graphicx}
\usepackage{xcolor}
\usepackage{amsmath,amssymb}
\usepackage{booktabs}
\usepackage{longtable}
\usepackage{array}
\usepackage{multirow}
\makeatletter
\@ifclassloaded{beamer}{%
  % enumitem 与 beamer 的 enumerate 模板冲突，幻灯片侧不加载
}{%
  \usepackage{enumitem}%
}
\makeatother
\usepackage{hyperref}
\usepackage{listings}
\usepackage{tcolorbox}
\tcbuselibrary{breakable,skins,listings}
\usepackage{tikz}
\usetikzlibrary{positioning,fit,arrows.meta,calc,matrix,backgrounds,shapes.geometric}

\definecolor{PressRed}{RGB}{160,50,40}
\definecolor{AdaptGreen}{RGB}{30,110,60}
\definecolor{PruneOrange}{RGB}{180,90,20}
\definecolor{PosBlue}{RGB}{30,80,140}
\definecolor{CodeBg}{RGB}{248,248,252}
\definecolor{CodeFrame}{RGB}{200,200,210}
\definecolor{CodeKeyword}{RGB}{0,0,180}
\definecolor{CodeComment}{RGB}{100,100,100}
\definecolor{CodeString}{RGB}{163,21,21}

\newtcolorbox{philobox}[1][]{
  colback=blue!4!white, colframe=PosBlue, fonttitle=\bfseries,
  title={#1}, breakable, arc=2pt
}
\newtcolorbox{pressbox}[1][]{
  colback=red!4!white, colframe=PressRed, fonttitle=\bfseries,
  title={#1}, breakable, arc=2pt
}
\newtcolorbox{adaptbox}[1][]{
  colback=green!4!white, colframe=AdaptGreen, fonttitle=\bfseries,
  title={#1}, breakable, arc=2pt
}
\newtcolorbox{prunebox}[1][]{
  colback=orange!5!white, colframe=PruneOrange, fonttitle=\bfseries,
  title={#1}, breakable, arc=2pt
}

\tikzset{
  evt/.style={rectangle, rounded corners=2pt, draw=gray!60, fill=blue!5,
    align=left, minimum width=9cm, minimum height=0.75cm, font=\small},
  pres/.style={rectangle, draw=PressRed!70, fill=red!6, align=center,
    minimum width=2.4cm, minimum height=0.65cm, font=\footnotesize},
  comp/.style={rectangle, draw=PosBlue!70, fill=blue!6, align=center,
    minimum width=2.6cm, minimum height=0.7cm, font=\footnotesize},
  arrow/.style={-{Stealth[length=2mm]}, thick, draw=gray!60}
}

% 仓库根相对本目录：../../../
\newcommand{\repopath}[1]{\filepath{../../../#1}}

% Workbench GUI 字体（gui/public/fonts）
\newcommand{\ebguifontpath}{../../../gui/public/fonts/}

\makeatletter
\@ifclassloaded{beamer}{\usefonttheme{professionalfonts}}{}
\makeatother

\setmainfont{JetBrainsMonoNLNerdFont-Regular}[
  Path=\ebguifontpath,
  Extension=.ttf,
  BoldFont=JetBrainsMonoNLNerdFont-Bold,
  ItalicFont=JetBrainsMonoNLNerdFont-Italic,
  BoldItalicFont=JetBrainsMonoNLNerdFont-BoldItalic,
  Scale=0.92
]
\setmonofont{JetBrainsMonoNLNerdFont-Regular}[
  Path=\ebguifontpath,
  Extension=.ttf,
  BoldFont=JetBrainsMonoNLNerdFont-Bold,
  ItalicFont=JetBrainsMonoNLNerdFont-Italic,
  BoldItalicFont=JetBrainsMonoNLNerdFont-BoldItalic
]
\newfontfamily\ebproggy{ProggyCleanNerdFontMono-Regular}[
  Path=\ebguifontpath,
  Extension=.ttf
]

% 中文回退：Windows 用雅黑（与 GUI --font-ui 一致），WSL/TeX Live 用 Fandol
\IfFontExistsTF{Microsoft YaHei}{%
  \setCJKmainfont{Microsoft YaHei}[
    BoldFont={Microsoft YaHei Bold},
    ItalicFont={Microsoft YaHei},
    BoldItalicFont={Microsoft YaHei Bold}
  ]
  \setCJKsansfont{Microsoft YaHei}[
    BoldFont={Microsoft YaHei Bold}
  ]
}{%
  \setCJKmainfont{FandolSong-Regular}[
    Path=/usr/share/texlive/texmf-dist/fonts/opentype/public/fandol/,
    Extension=.otf,
    BoldFont={FandolSong-Bold},
    ItalicFont={FandolSong-Regular},
    BoldItalicFont={FandolSong-Bold}
  ]
  \setCJKsansfont{FandolHei-Regular}[
    Path=/usr/share/texlive/texmf-dist/fonts/opentype/public/fandol/,
    Extension=.otf,
    BoldFont={FandolHei-Bold}
  ]
}

\lstdefinestyle{ebbackupcpp}{
  language=C++,
  basicstyle=\ebproggy\footnotesize,
  keywordstyle=\color{CodeKeyword}\bfseries,
  commentstyle=\color{CodeComment}\itshape,
  stringstyle=\color{CodeString},
  numbers=left,
  numberstyle=\tiny\color{gray},
  stepnumber=1,
  numbersep=6pt,
  backgroundcolor=\color{CodeBg},
  frame=single,
  rulecolor=\color{CodeFrame},
  breaklines=true,
  showstringspaces=false,
  tabsize=2,
  captionpos=b
}
\lstdefinestyle{ebbackupjson}{
  basicstyle=\ebproggy\footnotesize,
  numbers=left,
  numberstyle=\tiny\color{gray},
  backgroundcolor=\color{CodeBg},
  frame=single,
  rulecolor=\color{CodeFrame},
  breaklines=true,
  tabsize=2
}
\lstdefinestyle{ebbackupts}{
  language=Java,
  basicstyle=\ebproggy\footnotesize,
  keywordstyle=\color{CodeKeyword}\bfseries,
  commentstyle=\color{CodeComment}\itshape,
  stringstyle=\color{CodeString},
  numbers=left,
  numberstyle=\tiny\color{gray},
  backgroundcolor=\color{CodeBg},
  frame=single,
  rulecolor=\color{CodeFrame},
  breaklines=true,
  tabsize=2
}

% 答辩演讲稿共享宏
\newcommand{\vtag}[1]{\texttt{v#1}}
\newcommand{\gitdesc}[1]{\texttt{#1}}
\newcommand{\filepath}[1]{\texttt{\small #1}}
\newcommand{\wave}[1]{\textbf{Wave #1}}

% 逐字稿：双时长标注
\newcommand{\timing}[2]{%
  \par\smallskip\noindent{\small\color{gray}%
    \textbf{口播时长}\quad 课程版：#1 \quad|\quad 技术版：#2}\par\smallskip}

% 口播正文
\newenvironment{say}{%
  \begin{tcolorbox}[colback=blue!3!white,colframe=PosBlue!60,arc=2pt,boxrule=0.4pt,left=4pt,right=4pt,top=3pt,bottom=3pt]
  \small
}{%
  \end{tcolorbox}
}

% 演示动作提示
\newcommand{\demo}[1]{%
  \par\noindent{\small\color{AdaptGreen!80!black}\textbf{[演示]}\ #1}\par}

% Beamer 幻灯片标题副标
\newcommand{\slidesubtitle}{ebbackup 数据备份软件 · \vtag{0.9.5}}

% 代码片段（\lstinputlisting 包装）
\newcommand{\codesnippet}[4]{%
  \lstinputlisting[style=#1,firstline=#2,lastline=#3]{../../../#4}}

\newcommand{\showcifloor}{%
  \codesnippet{ebbackupjson}{1}{13}{engine_cpp/bench/baselines/ci_floor.json}}

\newcommand{\showrepostats}{%
  \codesnippet{ebbackupcpp}{38}{59}{engine_cpp/src/store/repo_stats.cc}}

\newcommand{\showrunbackup}{%
  \codesnippet{ebbackupcpp}{764}{768}{engine_cpp/src/engine/backup_engine.cc}}

\newcommand{\showshim}{%
  \codesnippet{ebbackupcpp}{124}{134}{engine_cpp/workbench/ebbackup_workbench_shim.cc}}

\newcommand{\showpanelsurface}{%
  \codesnippet{ebbackupts}{141}{149}{gui/src/utils/themeSettings.ts}}


\title{ebbackup 答辩演讲稿\\（课程版 8--10 分钟 + 技术版 15--20 分钟）}
\author{recoveryProjects / 《软件开发综合实验》}
\date{2026-07-07 \quad commit \gitdesc{1fca9d7}}

\begin{document}
\maketitle
\tableofcontents
\newpage

%% ============================================================
\part{Part I\quad 课程答辩口径（合计 8--10 分钟）}
%% ============================================================

\section{开场与项目定位}
\timing{0--1\,min}{0--1.5\,min}
\begin{say}
各位老师好。我是 ebbackup 数据备份软件项目的答辩人。本项目对应《软件开发综合实验》要求，实现了一款可在 Windows 与 Linux 上构建运行的\textbf{内容寻址增量备份}系统，版本为 \vtag{0.9.5}，包含 C++ 内核、命令行工具 \filepath{eb}、以及可选的桌面 Workbench 图形界面。

我们的目标不是做一个「能拷贝文件的演示程序」，而是在\textbf{可验证、可回归}的前提下，完成课程要求的数据备份与还原，并在压缩、加密、筛选、定时与实时备份等扩展项上给出可运行的工程实现。接下来我会按课程评分逻辑，先讲功能对照，再演示主流程，最后说明软件工程过程与测试门禁。
\end{say}

\section{课程功能对照}
\timing{1--2.5\,min}{1.5--3\,min}
% 课程基础/扩展功能对照（《软件开发综合实验》）
\begin{table}[htbp]
\centering
\caption{课程评分项与 ebbackup 实现对照}
\small
\begin{tabular}{@{}llp{5.8cm}@{}}
\toprule
课程项 & 分值 & 本项目实现 \\
\midrule
数据备份 & 基础40 & \filepath{eb backup} / C API / Workbench 备份页 \\
数据还原 & 基础40 & 全量/选择性/\texttt{--at} 时间旅行恢复 \\
特殊文件 & 10 & symlink/junction/目录树（\filepath{special\_files\_test}） \\
元数据 & 10 & mtime/atime/mode/uid/gid + \texttt{ApplyFileMeta} \\
自定义筛选 & 各3 & filter 文件：路径/glob/尺寸/时间/uid \\
打包解包 & 各10 & EbPack 内容寻址 pack（\filepath{.ebpack}） \\
压缩 & 各10 & LZ4 + zstd（ContentClass 自适应） \\
加密 & 各10 & AES-256-GCM + PBKDF2 \\
图形界面 & 10 & Tauri 2 + Vue Workbench（\filepath{gui/}） \\
定时备份 & 10 & \texttt{eb schedule} + GFS prune \\
实时备份 & 15 & \texttt{eb watch} + \texttt{FsWatch} \\
网络备份 & 10+ & \textbf{刻意不做}（生态边界） \\
\bottomrule
\end{tabular}
\end{table}

\begin{adaptbox}[答辩自述：第三方库规则]
课程规定扩展功能若用第三方库「直接实现」，扩展分记 50\%。LZ4/zstd/crypto 属成熟库；答辩重点放在 CDC、耐久性、Pipeline 与 bench 门禁等自研内核部分。
\end{adaptbox}

\begin{say}
课程基础分要求两项：将目录树备份到指定位置，以及从备份还原目录树。我们通过 \filepath{eb init}、\filepath{eb backup}、\filepath{eb restore} 三条命令完成，并在 Workbench 的六个 Activity 中提供同等能力。

扩展方面：我们实现了特殊文件与元数据、filter 筛选、EbPack 打包存储、LZ4 与 zstd 压缩、AES-GCM 加密、Tauri 桌面 GUI、\filepath{schedule} 定时与 \filepath{watch} 实时备份。网络网盘模式属于我们刻意划出的生态边界——单机可验证内核优先，因此未做用户管理与远程传输。

需要主动说明的是：LZ4、zstd 与加密库属于成熟第三方实现，按课程规则扩展分可能按 50\% 计；我们的答辩重点放在 CDC 分块、提交点耐久性、Pipeline 拓扑与 bench 门禁等\textbf{自研内核}部分。
\end{say}

\section{现场演示脚本（可录 3 分钟视频）}
\timing{2.5--5\,min}{3--6\,min}
\demo{启动 ebbackup Workbench；或终端执行 \texttt{eb} 命令}
% 备份主路径顺序图（简化）
\begin{figure}[htbp]
\centering
\begin{tikzpicture}[node distance=0.7cm]
  \node[evt, minimum width=11cm] (u) {用户：\texttt{eb backup <repo> <source>} 或 Workbench};
  \node[evt, below=of u, minimum width=11cm] (s) {Scan：遍历源目录 + filter 筛选};
  \node[evt, below=of s, minimum width=11cm] (c) {Chunk：FastCDC/HCRBO + Digest（SHA256）};
  \node[evt, below=of c, minimum width=11cm] (e) {Encode：LZ4/zstd + 可选 AES-GCM};
  \node[evt, below=of e, minimum width=11cm] (st) {Store：EbPack append + chunk.idx};
  \node[evt, below=of st, minimum width=11cm] (m) {Commit：manifest + superblock + RAR 审计};
  \draw[arrow] (u)--(s)--(c)--(e)--(st)--(m);
\end{tikzpicture}
\caption{全量/增量备份主路径（顺序图口语版）}
\label{fig:seq-backup}
\end{figure}

\begin{say}
演示分五步。第一步，在「仓库」页创建或打开本地 repo——空仓库在首次备份前没有 manifest，统计为零是正常状态，我们已在 \vtag{0.9.5} 修复了空仓库打开时的报错。第二步，在「备份」页选择\textbf{真实存在}的源目录；空仓库首次备份应使用全量模式。第三步，在「快照」页查看 txn 列表。第四步，在「恢复」页将数据还原到目标目录，可选指定 txn\_id 做时间旅行。第五步，展示 Workbench 的壁纸与卡片/输出面板透明度调节——这是产品化联调中真实遇到并解决的 WebView2 渲染问题。

若时间有限，口播可只演示：创建仓库 $\rightarrow$ 备份一个小文件夹 $\rightarrow$ 查看快照列表。
\end{say}

\section{软件工程过程（轻量 vs 课程模板）}
\timing{5--7\,min}{6--9\,min}
% 适应性工程环（摘自 ecosystem-evolution ch03）
\begin{figure}[htbp]
\centering
\begin{tikzpicture}[node distance=0.95cm]
  \node[evt] (b) {Bench/Profile 或 Bug 信号（P1--P5）};
  \node[evt, below=of b] (h) {假设：改 X 可缓解（不碰 Content ID 宪法）};
  \node[evt, below=of h] (i) {实现 + parity/powerfail 验证};
  \node[evt, below=of i] (g) {ctest + \filepath{ebbackup\_bench\_check}};
  \node[evt, below=of g] (d) {默认开启 / opt-in / 文档 CHANGELOG};
  \node[evt, below=of d] (p) {负向？剪枝回退};
  \draw[arrow] (b)--(h)--(i)--(g)--(d);
  \draw[arrow] (d)--(p);
  \draw[arrow] (p.west) -- ++(-1.2,0) |- (b.west);
\end{tikzpicture}
\caption{测量驱动的适应性工程环（替代重型瀑布评审链）}
\label{fig:adapt-loop}
\end{figure}

\begin{say}
课程讲授瀑布模型、面向对象方法与 UML 四图。我们的实际过程更接近\textbf{测量驱动的短环}：bench 或 bug 信号 $\rightarrow$ 假设 $\rightarrow$ parity 与 powerfail 验证 $\rightarrow$ ctest 与 bench\_check 门禁 $\rightarrow$ CHANGELOG 与文档归档。

这与课程并不矛盾：我们同样经历了需求边界（课程评分表）、设计（分层架构与构件划分）、实现、测试与发布。差异在于：内核性能类问题无法在一次设计评审中纸面定稿，必须用实测与黄金 parity 验证。事后我们编写了工程技术手册与生态实录，可映射为课程要求的类图、顺序图与构件图——StarUML 正式图可在答辩材料中作为附录提交。

变更管理方面：使用 git 版本控制，\filepath{CHANGELOG.md} 记录每个版本的 why，GitHub Actions 在 push 时运行 ctest。
\end{say}

\section{测试与质量门禁}
\timing{7--8.5\,min}{9--11\,min}
\begin{say}
测试分三层。语义层：FastCDC stream 与 slice 的 bit-identical parity，以及 powerfail 注入验证提交点。回归层：233 个 GoogleTest 用例，覆盖 pipeline、快照、加密、特殊文件等。性能层：L1 到 L7 bench，其中 reuse 不低于 90\%、pipeline ratio 不低于 0.90 为\textbf{immutable}，永不可为追分而降低。

下方是 CI 实际使用的 \filepath{ci\_floor.json} 节选，即我们的「完成定义」之一。
\end{say}
\showcifloor

\begin{say}
课程 Part 收尾：我们交付了可构建源码、CMake 构建脚本、CLI 与 Workbench 可执行程序、以及配套技术文档。谢谢各位老师，下面进入技术版延展，或直接进入问答。
\end{say}

\newpage
%% ============================================================
\part{Part II\quad 技术深度延展（+10 分钟，合计 15--20 分钟）}
%% ============================================================

\section{外部压力与本土化}
\timing{—}{11--13\,min}
\begin{say}
技术答辩从五条外部压力展开：P1 性能可比性——本地备份的硬通货是 MB/s；P2 I/O 与 fsync 成本；P3 正确性与并发；P4 运维可交付；P5 单人维护带宽。本土化含义是：Windows Release 为主环境、bench 为立法工具、流程极短、ABI 冻结避免频繁迁移。

\vtag{0.4.6} 时 L5 约 76 MB/s；四轮 sprint 后 \vtag{0.9.4} 达 170--172 MB/s，相对基准约翻倍。profile 显示单文件路径 chunk 阶段 CDC 约占 82\%，这是 Sprint 4 仍继续优化 CDC 的数据依据。
\end{say}
% 实测数据表（Release Windows, v0.9.4--0.9.5）
\begin{table}[htbp]
\centering
\caption{Stage 3.1 CI 门禁 vs \vtag{0.9.4} 实测}
\small
\begin{tabular}{@{}lrrr@{}}
\toprule
级别 & floor & 实测 & 余量 \\
\midrule
L5 256MB 单文件 & 105 MB/s & 170--172 & +63\% \\
L6 2GB & 112 MB/s & 186 & +66\% \\
L7 32$\times$32MB & 550 MB/s & 646--658 & +17\% \\
L3b streaming ratio & 0.90 & 1.01--1.05 & +12\% \\
L2 增量 reuse & 90\% & $\geq$90\% & immutable \\
\bottomrule
\end{tabular}
\end{table}

\begin{table}[htbp]
\centering
\caption{L5 演进与 Sprint 4 止损实验}
\small
\begin{tabular}{@{}lrrl@{}}
\toprule
节点 & L5 MB/s & 相对上版 & 备注 \\
\midrule
\vtag{0.4.6} & $\sim$76 & — & 基准 \\
\vtag{0.6.0} & $\sim$139 & +83\% & EbPack 默认化 \\
\vtag{0.9.0} & 142--149 & $\approx$0 & Wave 5 先修并发 \\
\vtag{0.9.2} & 154 & +3--8\% & StreamingChunkCpuPipeline \\
\vtag{0.9.3} & 159--163 & +3--6\% & digest\_base \\
\vtag{0.9.4} Phase B & 170--172 & +4--8\% & CDC seg1 bulk \\
Phase A opt-in & $\sim$105 & \textbf{-39\%} & \textbf{默认关闭} \\
\bottomrule
\end{tabular}
\end{table}

\begin{table}[htbp]
\centering
\caption{Stage 3.2 剪枝数据（真实实验）}
\small
\begin{tabular}{@{}lr@{}}
\toprule
指标 & 数值 \\
\midrule
Stage 3.2 L5 aspirational & 280 MB/s \\
\vtag{0.9.4} 实测 & 171 MB/s（仅为目标 61\%） \\
若 3.2 作 CI blocking & \textbf{100\% 失败} \\
决策 & 3.2 降为 nightly，3.1 作 merge 门禁 \\
\bottomrule
\end{tabular}
\end{table}


\section{真实实验 I：Wave 5 并发崩溃}
\timing{—}{13--14.5\,min}
\begin{pressbox}[\vtag{0.9.0} 实测记录]
多 worker Pipeline v4 触发 ChunkStore 全局索引锁下的 heap corruption（异常码 \texttt{0xC0000005}）以及 \texttt{ForEachRecord} 重入死锁。同期 L5 约 142--149 MB/s，L3b ratio 一度逼近 immutable 下限 0.90。
\end{pressbox}
\begin{say}
这是一次典型的 P3 压过 P1 的时刻。若选择继续 chase L5 分数，用户面对的是随机崩溃，生态信用归零。我们的决策是：shard 化索引、tombstone 独立锁序、atomic finalize，用两到三个 commit 先把正确性修到 ctest 全绿，再进入 v0.9.2 流式拓扑 sprint。

口播结论：\textbf{immutable floor 保护 ratio，但不保护 crash}；正确性优先是本土化路径里唯一理性的排序。
\end{say}

\section{真实实验 II：Stage 3.2 与负向剪枝}
\timing{—}{14.5--16\,min}
\begin{say}
\vtag{0.9.1} 曾讨论将 L5 floor 提升到 280 MB/s（Stage 3.2）。实测 \vtag{0.9.4} 为 171 MB/s，仅达目标的 61\%。若将 280 作为 merge 阻塞，CI 将 100\% 失败，迫使团队作弊降 floor 或放弃合并——这会破坏生态信用。

数据驱动的剪枝是：Stage 3.1（L5=105）blocking；Stage 3.2 aspirational 仅 nightly 报告。同类地，Phase 3C CDC carry-boundary 在 parity 失败后整支回退，收益 0 MB/s，不进入默认路径。
\end{say}

\section{真实实验 III：CDC Phase A 止损}
\timing{—}{16--17\,min}
\begin{say}
Sprint 4 实现两条 CDC 路径：Phase B 在 stream feed 上做 seg1 bulk scan，默认开启，L5 170--172 MB/s。Phase A 使用整文件 FastCdcSlice cuts，环境变量 \texttt{EBBACKUP\_CDC\_FAST\_SLICE=1}，实测仅约 105 MB/s，相对损失 39\%，甚至低于 \vtag{0.6.0} 的 139。

这是一次清晰的负向实验：L3b ratio 仍可能达标，但绝对吞吐崩溃。immutable ratio 门禁捕不到此类劣化，因此我们增加\textbf{默认路径 + 绝对 MB/s}的产品守卫，Phase A 保持 opt-in。答辩时我会把这作为「测量立法」优于「架构辩论」的例证。
\end{say}

\section{真实实验 IV：产品化联调（Workbench）}
\timing{—}{17--18.5\,min}
% 物理构件图（课程「构件图」答辩口径）
\begin{figure}[htbp]
\centering
\begin{tikzpicture}[node distance=1.2cm and 1.4cm]
  \node[comp] (cli) {CLI\\\filepath{eb}};
  \node[comp, right=of cli] (dll) {SHARED DLL\\\filepath{ebbackup\_workbench}};
  \node[comp, right=of dll] (gui) {Workbench\\Tauri+Vue};
  \node[comp, below=1.4cm of dll] (core) {静态库\\\filepath{ebbackup\_core}};
  \draw[arrow] (cli) -- (core);
  \draw[arrow] (dll) -- (core);
  \draw[arrow] (gui) -- node[above,font=\tiny]{JSON shim} (dll);
  \node[below=0.3cm of core, font=\footnotesize, align=center] {ChunkStore / Pipeline / Manifest / Snapshot};
\end{tikzpicture}
\caption{构件依赖：GUI 不直连内核，经 JSON shim 调用 C API}
\label{fig:components}
\end{figure}

\begin{say}
\vtag{0.9.5} 交付 Workbench 后，我们遇到三类真实问题。其一，新建仓库尚无 manifest，\texttt{repo\_info} 报错——根因是 \texttt{ComputeRepoStats} 无条件读取 manifest；修复后空仓库返回零统计。其二，GUI 透明度滑块无效——根因是壁纸模式下父层不透明且 WebView2 对嵌套 color-mix 支持不稳定；改为 JS 预计算 rgba。其三，备份报 source path not found——增加路径存在性校验与首次全量提示。

下面是空仓库统计修复、源路径校验与透明度预计算的三段代表性代码。
\end{say}
\showrepostats
\showrunbackup
\showpanelsurface

\section{架构与构件（技术版收尾）}
\timing{—}{18.5--19.5\,min}
% 内核八层架构（答辩精简）
\begin{figure}[htbp]
\centering
\begin{tikzpicture}[
  layer/.style={rectangle, draw=PosBlue!60, fill=blue!5, minimum width=11cm,
    minimum height=0.55cm, font=\small}
]
  \node[layer] (l1) {分块层：FastCDC + HCRBO（CFI 增量复用）};
  \node[layer, below=0.12cm of l1] (l2) {摘要层：SHA256 Content ID（Legacy/Standard 双栈）};
  \node[layer, below=0.12cm of l2] (l3) {压缩层：ContentClass $\rightarrow$ LZ4/zstd/跳过};
  \node[layer, below=0.12cm of l3] (l4) {存储层：ChunkStore + EbPack 16-shard};
  \node[layer, below=0.12cm of l4] (l5) {元数据：Manifest v4 + 双槽 SuperBlock};
  \node[layer, below=0.12cm of l5] (l6) {时间旅行：Snapshot + GFS prune};
  \node[layer, below=0.12cm of l6] (l7) {可验证：RAR 审计链 + Merkle};
  \node[layer, below=0.12cm of l7] (l8) {并发：Pipeline v4 / StreamingChunkCpuPipeline};
\end{tikzpicture}
\caption{ebbackup 内核分层（\filepath{kernel-manual} 概述）}
\label{fig:arch-layers}
\end{figure}

\begin{philobox}[生态宪法（永不交易）]
\begin{itemize}
  \item Content ID = \texttt{SHA256(content)}，CDC 路径不得改变语义；
  \item 提交点耐久性：manifest commit 前 \texttt{Flush}/fsync；
  \item Immutable CI：\texttt{reuse\_pct$\geq$90}、\texttt{pipeline\_ratio*$\geq$0.90} 永不可降；
  \item 中断 txn $\rightarrow$ \texttt{kAborted}，前一 snapshot 仍可 verify。
\end{itemize}
\end{philobox}

\showshim
\begin{say}
技术版收尾：内核八层、构件图与 JSON shim 边界。DLL 不复制业务逻辑，只做 C API 的 JSON 封装，便于 Tauri Rust 侧集成测试。完整 API 见 \filepath{kernel-manual} 第 18 章，生态决策见 \filepath{ecosystem-evolution}。
\end{say}

\newpage
%% ============================================================
\part{附录\quad Q\&A 备答（每题约 30 秒）}
%% ============================================================

\section{常见问题}
\begin{enumerate}
  \item \textbf{EbPack 算不算课程「打包解包」？}\\
  答：EbPack 将 chunk 聚合为内容寻址的 \filepath{.ebpack} 文件，是备份领域的 pack 存储，而非 tar 式简单拼接；答辩时强调 dedup 与索引，而非「单文件归档」。

  \item \textbf{为何不用 Qt？}\\
  答：课程允许 UI 用脚本技术栈；我们选用 Tauri+Vue 以复用 eB-Tree Workbench 壳层经验，内核仍为 C++，符合「后台 C++、前台可脚本」规则。

  \item \textbf{为何不做网络备份？}\\
  答：生态边界是单节点可验证内核；网络协议、用户管理与传输加密会分散维护带宽，与 P5 压力冲突。

  \item \textbf{与 restic/borg 的差异？}\\
  答：同为主机级内容寻址备份；我们强调 Windows Release bench 门禁、EbPack 默认布局、RAR/Merkle 审计链与 C API 嵌入场景。

  \item \textbf{immutable floor 是什么？}\\
  答：CI 中 reuse 与 pipeline ratio 下限只允许升高不允许降低，防止为追分牺牲增量正确性。

  \item \textbf{单人项目如何替代 Scrum？}\\
  答：用 ctest+bench\_check 作为机器可执行的「完成定义」；CHANGELOG 替代 Sprint 计划文档。

  \item \textbf{UML 四图在哪里？}\\
  答：内核手册与本文档含类/顺序/构件的 tex 表述；StarUML 正式图可作为课程提交附录，从事后架构导出。

  \item \textbf{下一步？}\\
  答：Sprint 5 候选是 Hybrid ChunkCuts+feed overlap；Stage 3.2 的 280 MB/s 需再 +109 MB/s，以数据而非口号推进。
\end{enumerate}

\end{document}
